\section{The \texttt{Pol\_mirror} McStas Component}
Polarising mirror.

\subsection*{Identification}
\begin{itemize}
  \item \textbf{Site:} 
  \item \textbf{Author:} Peter Christiansen
  \item \textbf{Origin:} RISOE
  \item \textbf{Date:} July 2006
\end{itemize}

\subsection*{Description}
\begin{lstlisting}
This component models a rectangular infinitely thin mirror.
For an unrotated component, the mirror surface lies in the Y-Z
plane (ie. parallel to the beam).
It relies on similar physics as the Monochromator_pol.
The reflectivity function (see e.g. share/ref-lib for examples) and parameters
are passed to this component to give a bigger freedom.
The up direction is hardcoded to be along the y-axis (0, 1, 0)
IMPORTANT: At present the component only works correctly for polarization along the up/down
direction and for completely unpolarized beams, i.e. sx=sy=sz=0 for all rays.

For now we assume:
P(Transmit|Q) = 1 - P(Reflect|Q)
i.e. NO ABSORPTION!


The component can both reflect and transmit neutrons with a respective proportion
depending on the p_reflect parameter:
p_reflect=-1 Reflect and transmit (proportions given from reflectivity) [default]
p_reflect=1 Only handle reflected events
p_reflect=0 Only handle transmitted events (reduce weight)
p_reflect=0-1 Both transmit and reflect with fixed statistics proportions

The parameters can either be
double pointer initializations (e.g. {R0, Qc, alpha, m, W})
or table names (e.g."supermirror_m2.rfl" AND useTables=1).
NB! This might cause warnings by the compiler that can be ignored.

Examples:
Reflection function parametrization
Pol_mirror(zwidth = 0.40, yheight = 0.40, rUpFunc=StdReflecFunc, rUpPar={1.0, 0.0219, 6.07, 2.0, 0.003})

Table function
Pol_mirror(zwidth = 0.40, yheight = 0.40, rUpFunc=TableReflecFunc, rUpPar="supermirror_m2.rfl", rDownFunc=TableReflecFunc, rDownPar="supermirror_m3.rfl", useTables=1)

See also the example instrument Test_Pol_Mirror (under tests).

GRAVITY: YES

%BUGS
NO ABSORPTION
\end{lstlisting}

\subsection*{Input parameters}
Parameters in \textbf{boldface} are required; the others are optional.

\begin{longtable}{p{0.22\textwidth}p{0.12\textwidth}p{0.46\textwidth}p{0.14\textwidth}}
\toprule
\textbf{Name} & \textbf{Unit} & \textbf{Description} & \textbf{Default} \\
\midrule
\endhead
rUpPar & 1 & Parameters for rUpFunc & \{0.99,0.0219,6.07,2.0,0.003\} \\
rDownPar & 1 & Parameters for rDownFunc & \{0.99,0.0219,6.07,2.0,0.003\} \\
rUpData & str & Mirror Reflectivity data file for spin up & "" \\
rDownData & str & Mirror Reflectivity data file for spin down & "" \\
p\_reflect & 1 & Proportion of reflected events. Use 0 to only get the transmitted beam, and 1 to get only the reflected beam. Use -1 to use the mirror reflectivity. This value is purely computational and is not related to the actual reflectivity & -1 \\
\textbf{zwidth} & m & Width of the mirror &  \\
\bottomrule
\end{longtable}

\subsection*{Links}
\begin{itemize}
  \item \href{run:/home/willend/willend-McCode/mcstas-comps/optics/Pol_mirror.comp}{Source code} for \texttt{Pol\_mirror.comp}.
\end{itemize}
\IfFileExists{Pol_mirror_static.tex}{\input{Pol_mirror_static.tex}}{}